\chapter{Conclusiones}
\label{cap:conclusiones}

% TODO: redactar las conclusiones finales una vez integrados:
%   - Capitulo 4 (SOAP vs REST, Carvajal)
%   - Capitulo 5 (Tendencias web, Fajardo)
%   - Capitulo 6 (Revision cruzada consolidada, Carvajal + Fajardo)
%
% Las conclusiones deben derivarse de esos aportes y de los resultados del
% capitulo 7, no redactarse de forma independiente ni anticipada.

\pendienteIntegracion{Conclusiones finales del informe de la GA de Unidad IV. Se redactar\'an a partir de los cap\'itulos~\ref{cap:soap-vs-rest}, \ref{cap:tendencias-web}, \ref{cap:revision-cruzada}\ y~\ref{cap:resultados}, una vez integrados los aportes de Carvajal Loor Johan Stalin y Fajardo Montes Michael Xavier.}

\section{Alcance previsto}
\label{sec:conclusiones-alcance}

Las conclusiones finales de este informe abordar\'an, como m\'inimo:

\begin{itemize}[nosep]
  \item El cumplimiento t\'ecnico de BIOPET frente al alcance definido
    para esta actividad, con base en la evidencia del
    cap\'itulo~\ref{cap:resultados}.
  \item Las principales fortalezas y debilidades identificadas en la
    revisi\'on cruzada consolidada (cap\'itulo~\ref{cap:revision-cruzada}).
  \item Las lecciones aprendidas de la comparaci\'on SOAP vs REST aplicadas
    al caso real de BIOPET (cap\'itulo~\ref{cap:soap-vs-rest}).
  \item La pertinencia de las tendencias web analizadas
    (cap\'itulo~\ref{cap:tendencias-web}) para la evoluci\'on futura del
    sistema.
\end{itemize}

No se redactan conclusiones definitivas en esta versi\'on del informe
porque dependen directamente de contenido todav\'ia no integrado.
